% V18 integration of Dot batch 099; immutable source SHA256 85f98276379eff53d810905d660c7797ab074ad35c8064a8fa64c28c8a805536
\begingroup
% These definitions have only group-local scope; the shared preamble is untouched.
\def\OO{\mathcal O}\def\FF{\mathcal F}\def\GG{\mathcal G}
\def\PP{\mathbf P}\def\calP{\mathcal P}\def\dd{\delta}
\def\Div{\operatorname{Div}}\def\depth{\operatorname{depth}}
\def\pr{\operatorname{pr}}\def\coker{\operatorname{coker}}
\egas{7.4}{命题：译校采用的附加条件版本}
假定 $g:X\to S$ 射影、平坦、有限表示，取 $s\in S$，并令 $n=\dim X_s$。再假定 $X_s$ 非空且纯维 $n$，且 $X_s$ 的非光滑轨迹维数小于 $n$。
则存在 $s$ 的开邻域 $U$ 及有限、平展、满射的 $S'\to U$，使 $X\times_S S'\to S'$ 有截面；
若 $k(s)$ 无限，还可取 $S'$ 为 $X_U$ 的闭子概形。\cmt{D18-099-7.4}{
这是第99批译校补证中的射影、有限表示、纯维版本，不是恢复出的法文原命题。
法文物理第29页原写固有、平坦，并要求不可分点集的维数严格小于 $n$；英文物理第20页误作 $\le n$。
纯不可分反例否定的是英译的放宽，不能据此否定法稿的严格不等号版本。
法稿一般陈述及原归纳路线完整见本版核查报告“D18-099：（7.4）一般原陈述、证明路线及限制”；
在一般固有情形，所给切割证明仍待核。}
\begin{proof}[译校补证，非原文]
选取嵌入$X\hookrightarrow\mathbf P^N_S$。在几何纤维上，余维 $n$ 的线性子空间中，与非光滑轨迹相交者构成 Grassmann 丛参数空间的真闭子集：对固定点，包含该点施加 $n$ 个条件，故上述轨迹的维数不足以覆盖参数空间。在光滑点上，包含该点的子空间若还不横截切空间，再多至少一个独立的闭条件。将该关联坏轨迹投到参数空间，其闭包维数仍严格小于参数空间。因而存在非空开集，其任意参数给出非空、有限、几何约化的横截；非空由射影维数不等式保证。这也解释了英46页作者提出直接取余维 $n$ 的截面而不用逐次切割。

若 $k(s)$ 无限，利用 Grassmann 丛的仿射大图可在该开集中选 $k(s)$-点，并在 $s$ 的开邻域提升矩阵坐标。所得闭截面在纤维上的各点处，由底上的光滑性及方程微分线性无关，对底 étale。该截面相对底固有，所以可缩小底去除非 étale 点的闭像。由固有且拟有限得有限，由非空纤维与开映射得缩小后的满射。取此截面自身为 $S'$，闭嵌入给出所需的换底截面。若 $k(s)$ 有限，可先取好参数开集中的一个闭点；其剩余域有限可分。按7.1作纤维只有该点的有限 étale 换底，提升其参数。若需在新底缩邻域，可先从原底删去新底坏闭集在有限态射下的闭像；最终有限 étale 满射可与第一次换底复合。
\end{proof}

\egas{7.5}{备注（关于射影性）}
作者指出，射影且相对可分的族，在所讨论的非空纤维附近可有 étale 多截面。作者接着声称，可以给出固有、光滑但非射影而使同样结论失败的例子；他还说，一般不能把射影假设减弱为拟射影假设，并提出取非有限的 étale 态射作为反例的方向。\footnote{法29这一段只写明寻找反例的方向；页末手记的末尾读法未定。没有给出一个完整的固有、光滑而非射影的构造。非有限的 étale 态射本身不足以说明无有限 étale 多截面。若仍同时要求固有，则拟射影态射已经是射影的。}

\par\noindent\textit{译校补证另见核查报告。}\cmt{D18-099-R}{英译不等号反例、零维基例、较弱的局部结论及拟射影反例完整保存在本版核查报告“D18-099”条；其中译校补证为非原文。}

\endgroup
